% Part: proof-theory
% Chapter: normalization
% Section: segments

\documentclass[../../../include/open-logic-section]{subfiles}

\begin{document}

\olfileid{pt}{nor}{seg}

\olsection{المقاطع والقطوع}

قاعدة~!!^a{introduction} تتبعها قاعدة~!!a{elimination} هي التفاف في
!!a{proof}، وتطبيع~!!a{proof} هو حذف مثل هذه الالتفافات. غير أن قاعدتي
\Elim{\lor} و\Elim{\lexists} تجعلان حتى تعريف ``الالتفاف الذي نريد حذفه''
معقدًا بعض الشيء. والمشكلة هي أن هاتين القاعدتين لا تضمنان أن قاعدة
!!{elimination} تتبع قاعدة~!!{introduction} \emph{مباشرةً} داخل التفاف؛
فقد تقع~\Elim{\lor} أو~\Elim{\lexists} بينهما، مثلًا:
\[
\AxiomC{}
\DeduceC{$\lexists[x][!C(x)]$}
\AxiomC{$\Discharge{!A}{x}$}
\DeduceC{$!B$}
\DischargeRule{\Intro{\lif}}{x}
\UnaryInfC{$!A \lif !B$}
\DischargeRule{\Elim{\lexists}}{y}
\BinaryInfC{$!A \lif !B$}
\AxiomC{}
\DeduceC{$!A$}
\RightLabel{\Elim{\lif}}
\BinaryInfC{$!B$}
\DisplayProof
\]
وفي الواقع قد يفصل أي عدد من تطبيقات~\Elim{\lor} أو~\Elim{\lexists}
بين~\Intro{\lif} و\Elim{\lif}، مثلًا:
\[
\AxiomC{}
\DeduceC{$\lexists[x][!C(x)]$}
\AxiomC{}
\DeduceC{$!D \lor !E$}
\AxiomC{$\Discharge{!A}{x}$}
\DeduceC{$!B$}
\DischargeRule{\Intro{\lif}}{x}
\UnaryInfC{$!A \lif !B$}
\AxiomC{$\Discharge{!A}{x}$}
\DeduceC{$!B$}
\DischargeRule{\Intro{\lif}}{x}
\UnaryInfC{$!A \lif !B$}
\RightLabel{\Elim{\lor}}
\TrinaryInfC{$!A \lif !B$}
\DischargeRule{\Elim{\lexists}}{y}
\BinaryInfC{$!A \lif !B$}
\AxiomC{}
\DeduceC{$!A$}
\RightLabel{\Elim{\lif}}
\BinaryInfC{$!B$}
\DisplayProof
\]
يبين هذا المثال أن تطبيق~\Elim{\lif} نفسه قد يدخل في التفافات متعددة، إذ
يتبع أكثر من تطبيق واحد لـ~\Intro{\lif}. ولشمول جميع هذه الاحتمالات، نعرّف
الالتفافات باستعمال أنواع معينة من متتاليات وقائع~!!{formula} في برهان
\Log{N1i}، ونسميها \emph{مقاطع}. وفي مثالنا هي متتالية من وقائع
$!A \lif !B$ تتبع فيها نتيجة~$\Elim{\lor}$ أو~$\Elim{\lexists}$ مقدمتها
الصغرى. وفي المثال مقطعان من هذا النوع، يحتوي كل منهما ثلاثة وقائع من
$!A \lif !B$. وهذا هو التعريف الرسمي.

\begin{defn}
\emph{المقطع} في~!!{proof}~$\delta$ من~\Log{N1i} هو متتالية الوقائع
$!A_1$، \dots، $!A_n$ لـ~!!{formula}~$!A$ نفسها في~$\delta$، بحيث:
\begin{enumerate}
  \item لا تكون~$!A_1$ نتيجة~\Elim{\lor} أو~\Elim{\lexists}؛
  \item لا تكون~$!A_n$ مقدمةً صغرى لـ~\Elim{\lor} أو~\Elim{\lexists}؛
  \item لكل~$i = 1$، \dots، $n-1$، تكون~$!A_i$ مقدمةً صغرى لـ
  \Elim{\lor} أو~\Elim{\lexists}، وتكون~$!A_{i+1}$ نتيجة ذلك الاستدلال.
\end{enumerate}
و\emph{طول} هذا المقطع هو~$n$.
\end{defn}

ليس كل مقطع هدفًا للحذف (أي التفافًا). وتسمى المقاطع التي تكون أهدافًا
\emph{قطوعًا}.

\begin{defn}
\emph{القطع} مقطع تكون فيه~$!A_n$ المقدمة الكبرى لاستدلال~!!a{elimination}،
ويكون إما~$n=1$ و$!A_1=!A_n$ نتيجة استدلال~!!a{introduction} أو~\FalseInt،
وإما~$n>1$.
\end{defn}

لاحظ أن مقطع القطع لا يلزم أن يبدأ بنتيجة قاعدة~!!a{introduction}. وكل ما
يلزم هو أن ينتهي بالمقدمة الكبرى لقاعدة~!!a{elimination}. لكن المقطع الذي
طوله~$1$ يجب أن يكون نتيجة قاعدة~!!{introduction} أو~\FalseInt كي يُعد
قطعًا.

\begin{defn}
\emph{رتبة} القطع هي~$\depth{!A}$. و\emph{رتبة القطع}
$\cutrank{\delta}$ لـ~!!a{proof}~$\delta$ هي أكبر رتب القطوع في~$\delta$،
وتُعرَّف بأنها~$0$ إذا لم توجد قطوع. ويكون القطع \emph{أقصى} في~$\delta$
إذا كانت رتبته~$\cutrank{\delta}$؛ أي إذا كانت رتبته أكبر ما يمكن.
و\emph{طول القطع} لـ~$\delta$ هو مجموع أطوال جميع قطوعه القصوى، ويُعرَّف
بأنه~$0$ إذا لم توجد قطوع.
\end{defn}

\end{document}
